\documentclass[a4paper,11pt]{jlreq}

\usepackage{amsmath,amssymb,bm}
\usepackage{booktabs}
\usepackage{enumitem}
\usepackage{graphicx}
\usepackage{hyperref}
\usepackage{longtable}
\usepackage[haranoaji]{luatexja-preset}
\usepackage{siunitx}
\usepackage{xcolor}

\hypersetup{
  colorlinks=true,
  linkcolor=blue!50!black,
  urlcolor=blue!50!black,
  citecolor=blue!50!black,
}
\setlist[itemize]{leftmargin=2em}
\setlist[enumerate]{leftmargin=2.2em}
\sisetup{
  group-separator = {,},
  group-minimum-digits = 4,
}
\graphicspath{
  {../../data/derived/m_trade_quote_cpcf_examples/M_20251031-20251103_min100_20260408T060741Z/}
  {../../data/derived/m_trade_quote_round_lot_effects/M_20251031-20251103_min100_20260408T060741Z_QZ_with_round_lot_sep2025_mean_close/}
  {../../data/derived/m_trade_quote_rd_250/M_20251031-20251103_min100_20260408T060741Z_QZ_with_round_lot_sep2025_mean_close/}
}

\title{M 接頭辞銘柄に対する trade-quote panel 構築と round-lot 分析メモ}
\author{2026-04-08 共有用}
\date{}

\begin{document}
\maketitle

\section{現状のデータ処理}

\subsection{条件}

\begin{itemize}
  \item 対象日: \(d \in \mathcal{D}=\{\texttt{20251031}, \texttt{20251103}\}\)
  \item 対象銘柄集合: 各日における Daily TAQ MASTER のうち，
        ticker が \texttt{M} で始まる銘柄
  \item 集計対象の取引所ペア: trade 側 \(e=\texttt{Q}\)，quote 側 \(g=\texttt{Z}\)
  \item 観測窓: \(\mathcal{W}=[35100.0,56700.0]\)，
        すなわち \texttt{09:45:00} から \texttt{15:45:00}
  \item 最低イベント数: 各 side で 100
  \item mode 推定の探索範囲: \(\mathcal{U}=[-1,1]\)
  \item bandwidth 候補:
        \(\mathcal{B}=\{10^{-7},10^{-6},10^{-5},10^{-4}\}\)
  \item kernel: tent
  \item simple 条件: trade 側と quote 側で event time が一致する pair は除外
\end{itemize}

以下で使う pre/post の日付はそれぞれ
\(\texttt{20251031}\) と \(\texttt{20251103}\) である．

\subsection{入力データ}

今回の panel 構築で用いる入力は 4 種類である．

\begin{center}
\begin{tabular}{p{0.18\linewidth}p{0.28\linewidth}p{0.42\linewidth}}
\toprule
入力 & 位置 & 用途 \\
\midrule
MASTER &
\texttt{data/dailyTAQ/MASTER/} &
M 接頭辞銘柄集合の抽出と，各日・各銘柄の \texttt{Round\_Lot} 参照に使う． \\
TRADE &
\texttt{data/dailyTAQ/TRADE/} &
trade 側イベント列の構成に使う． \\
BBO &
\texttt{data/dailyTAQ/BBO/} &
quote 側イベント列の構成に使う． \\
日次終値 &
\texttt{yfinance} &
\(\texttt{2025-09-01}\) から \(\texttt{2025-10-01}\) の
日次終値平均 \(\bar{p}^{\mathrm{close}}_s\) を構成する． \\
\bottomrule
\end{tabular}
\end{center}

各日 \(d\) の MASTER から，
\[
  \mathcal{S}_d
  =
  \left\{
    s
    \mid
    s \text{ starts with \texttt{M}}
  \right\}
\]
を作る．既存成果物では各日とも \(|\mathcal{S}_d|=636\) である．

\subsection{パイプライン全体像}

panel 構築の流れは次の 4 段階である．

\begin{enumerate}
  \item MASTER から M 接頭辞銘柄集合を抽出する．
  \item raw TRADE と raw BBO を観測窓 \(\mathcal{W}\) に制限し，
        銘柄別 filtered parquet を作る．
  \item 各 \((d,s)\) について取引所 pair を走査し，
        mode と cross-K を集計する．
  \item \(e=\texttt{Q}\)，\(g=\texttt{Z}\)，status=\texttt{ok} の pair だけを残し，
        2 日比較の symbol-level panel に整形する．
\end{enumerate}

このとき，最終的に使う panel を
\[
  P=\{p_s\}_{s=1}^{N_P}
\]
と書く．ここで \(s\) は銘柄 index であり，
\(N_P\) は panel に含まれる銘柄数である．

\subsection{最終 panel の構成}

最終 panel は，各銘柄 \(s\) に対して
\[
  p_s =
  \left(
    s,
    \bar{p}^{\mathrm{close}}_s,
    p_{s,\texttt{20251031}},
    p_{s,\texttt{20251103}}
  \right)
\]
を 1 行に持つ wide 形式の表である．
各日ブロック \(p_{s,d}\) には round lot，closest mode to zero，
6 個の cross-K 指標が入る．

列数は固定で 18 であり，今回の panel の行列サイズは \((153,18)\) である．
2 日とも値が揃う銘柄数は 130，そのうち
round lot が日付間で変化する銘柄数は 9 である．

\subsection{CPCF の具体例}

Section 2 と Section 3 の推定結果は，
最終的には各銘柄・各日について選ばれた bandwidth のもとでの
CPCF の形状に支えられている．
そこで，Q-Z pair の event-time 構造が実際にどのように見えるかを，
4 銘柄の具体例で示す．
横軸は trade 側と quote 側のイベント時刻のずれをマイクロ秒単位で表し，
縦軸はそのずれに対応する CPCF の値である．

選んだ 4 銘柄は次のとおりである．
\begin{itemize}
  \item \texttt{MORN}: \(250\) ドル cutoff の直下にあり，round lot が変化しない例
  \item \texttt{MAR}: \(250\) ドル cutoff の直上にあり，\(100 \to 40\) へ移行した例
  \item \texttt{MDB}: 40 へ移行した別の例
  \item \texttt{MSFT}: 高流動性の 40 への移行例
\end{itemize}

\begin{figure}[p]
\centering
\includegraphics[width=0.48\linewidth]{cpcf__MORN.png}
\includegraphics[width=0.48\linewidth]{cpcf__MAR.png}

\vspace{0.8em}

\includegraphics[width=0.48\linewidth]{cpcf__MDB.png}
\includegraphics[width=0.48\linewidth]{cpcf__MSFT.png}
\caption{CPCF の具体例．各図は \texttt{2025-10-31} と \texttt{2025-11-03} を左右に並べ，
各銘柄・各日に対して選ばれた bandwidth をそのまま使って CPCF を再評価したものである．
左上が \texttt{MORN}，右上が \texttt{MAR}，左下が \texttt{MDB}，右下が \texttt{MSFT} である．
太い縦線は 0 に最も近い mode 候補，細い縦線は検出された mode 候補を表す．
この 4 例は，Section 2 と Section 3 で比較している round-lot 関連の結果が，
日ごと・銘柄ごとの event-time 構造の違いの上に成り立っていることを示している．}
\end{figure}

\section{round lot 変更回帰と箱ひげ図}

\subsection{設計}

この節では，事前 round lot が 100 である銘柄のうち，
事後 round lot が 100 のまま残った群と 40 へ移行した群を比較する．
比較サンプルを
\[
  \mathcal{S}^{\mathrm{cmp}}
  =
  \left\{
    s
    \mid
    r^{\mathrm{pre}}_s = 100,\
    r^{\mathrm{post}}_s \in \{100,40\},\
    \bar{p}^{\mathrm{close}}_s \text{ が観測されている}
  \right\}
\]
と書く．ここで \(r^{\mathrm{pre}}_s\) は pre round lot，
\(r^{\mathrm{post}}_s\) は post round lot である．

40 への移行指標を
\[
  D_s = \mathbf{1}\{r^{\mathrm{post}}_s = 40\}
\]
と定義する．
アウトカムは \(Y_s\) と書き，その変化を \(\Delta Y_s\) とする．
\(\Delta Y_s\) は \texttt{mode} では事後値から事前値を引いた単純差，
cross-K では \(\log(1+Y^{\mathrm{post}}_s) - \log(1+Y^{\mathrm{pre}}_s)\)
とする．

比較回帰は二群 OLS
\[
  \Delta Y_s = \alpha + \beta D_s + \varepsilon_s
\]
に固定する．ここで \(\beta\) は
40 へ移行した群の平均変化から
100 のままの群の平均変化を引いた差を表す．

\subsection{サンプル構成}

今回の分析結果では，サンプル構成は次のとおりである．

\begin{center}
\begin{tabular}{lr}
\toprule
項目 & 値 \\
\midrule
最終 panel の総行数 & 153 \\
比較サンプルの行数 & 121 \\
40 へ移行した群 & 8 \\
100 のままの群 & 113 \\
事後 round lot \(=10\) のため除外 & 1 \\
\(250\) ドル基準と整合的な割合（100 のままの群） & 0.973451 \\
\(250\) ドル基準と整合的な割合（40 へ移行した群） & 1.000000 \\
\bottomrule
\end{tabular}
\end{center}

\subsection{回帰結果の要約}

7 つのアウトカムそれぞれについて，
40 へ移行した群の平均変化から
100 のままの群の平均変化を引いた値を推定すると，次の結果になる．

\begin{table}[htbp]
\centering
\small
\resizebox{\linewidth}{!}{
\begin{tabular}{l l r r r r r r}
\toprule
アウトカム & 変化の定義 & 平均変化の差 & 95\% CI 下限 & 95\% CI 上限 & p値 & 40 へ移行 & 100 のまま \\
\midrule
\texttt{cross\_k\_neg\_1e3\_1e4} & \( \log(1+Y^{\mathrm{post}})-\log(1+Y^{\mathrm{pre}}) \) &  0.001660 & -0.129082 &  0.132403 & 0.980144 & 8 & 113 \\
\texttt{cross\_k\_neg\_1e4\_1e5} & \( \log(1+Y^{\mathrm{post}})-\log(1+Y^{\mathrm{pre}}) \) & -0.009402 & -0.086283 &  0.067479 & 0.810569 & 8 & 113 \\
\texttt{cross\_k\_neg\_1e5\_0}   & \( \log(1+Y^{\mathrm{post}})-\log(1+Y^{\mathrm{pre}}) \) & -0.007898 & -0.030912 &  0.015116 & 0.501183 & 8 & 113 \\
\texttt{cross\_k\_pos\_0\_1e5}   & \( \log(1+Y^{\mathrm{post}})-\log(1+Y^{\mathrm{pre}}) \) & -0.017276 & -0.040729 &  0.006177 & 0.148811 & 8 & 113 \\
\texttt{cross\_k\_pos\_1e4\_1e3} & \( \log(1+Y^{\mathrm{post}})-\log(1+Y^{\mathrm{pre}}) \) & -0.098694 & -0.285409 &  0.088021 & 0.300202 & 8 & 113 \\
\texttt{cross\_k\_pos\_1e5\_1e4} & \( \log(1+Y^{\mathrm{post}})-\log(1+Y^{\mathrm{pre}}) \) & -0.033201 & -0.109971 &  0.043569 & 0.396641 & 8 & 113 \\
\texttt{mode}                    & \( Y^{\mathrm{post}}-Y^{\mathrm{pre}} \)                & -0.000006 & -0.000017 &  0.000006 & 0.331071 & 8 & 113 \\
\bottomrule
\end{tabular}
}
\caption{round-lot change 比較回帰の要約．}
\end{table}

5\% 水準で有意な群間差は確認されない．
係数の絶対値が最も大きいのは
\texttt{cross\_k\_pos\_1e4\_1e3} であり，
次が \texttt{cross\_k\_pos\_1e5\_1e4} と
\texttt{cross\_k\_pos\_0\_1e5} である．
ただし，いずれも信頼区間は 0 を跨いでいる．

\subsection{係数プロット}

\begin{figure}[htbp]
\centering
\includegraphics[width=0.86\linewidth]{coef_plot.png}
\caption{各アウトカムについて，
40 へ移行した群の平均変化から
100 のままの群の平均変化を引いた値と 95\% 信頼区間．}
\end{figure}

\subsection{箱ひげ図}

\begin{figure}[p]
\centering
\includegraphics[width=0.78\linewidth]{boxplot__mode.png}
\caption{\texttt{mode} の箱ひげ図．図中の 2 群は，
事後 round lot が 100 のままの群と 40 へ移行した群を表す．}
\end{figure}

\begin{figure}[p]
\centering
\includegraphics[width=0.32\linewidth]{boxplot__cross_k_neg_1e3_1e4.png}
\includegraphics[width=0.32\linewidth]{boxplot__cross_k_neg_1e4_1e5.png}
\includegraphics[width=0.32\linewidth]{boxplot__cross_k_neg_1e5_0.png}
\caption{負側の cross-K 変化の箱ひげ図．左から
\texttt{cross\_k\_neg\_1e3\_1e4},
\texttt{cross\_k\_neg\_1e4\_1e5},
\texttt{cross\_k\_neg\_1e5\_0}.}
\end{figure}

\begin{figure}[p]
\centering
\includegraphics[width=0.32\linewidth]{boxplot__cross_k_pos_0_1e5.png}
\includegraphics[width=0.32\linewidth]{boxplot__cross_k_pos_1e5_1e4.png}
\includegraphics[width=0.32\linewidth]{boxplot__cross_k_pos_1e4_1e3.png}
\caption{正側の cross-K 変化の箱ひげ図．左から
\texttt{cross\_k\_pos\_0\_1e5},
\texttt{cross\_k\_pos\_1e5\_1e4},
\texttt{cross\_k\_pos\_1e4\_1e3}.}
\end{figure}

\subsection{解釈}

この節の回帰は cutoff を使う RD ではなく，
実際に 40 へ移行した 8 銘柄と
100 のまま残った 113 銘柄の比較である．
両群の価格帯は大きく異なり，
平均価格は 100 のままの群がおよそ 72.4，
40 へ移行した群がおよそ 497.9 である．
したがって，ここでの \(\beta\) は round-lot 変更の因果効果というより，
移行群と非移行群の記述的比較として読むべきである．

\section{\$250\$ cutoff の RD 分析}

\subsection{RD 設計}

RD では走行変数を
\[
  X_s = \bar{p}^{\mathrm{close}}_s - 250
\]
と定義し，cutoff を上回るかどうかの指標を
\[
  D^{\mathrm{rule}}_s = \mathbf{1}\{X_s > 0\}
\]
と置く．
RD サンプルでは \(\bar{p}^{\mathrm{close}}_s < 1000\) を課し，
さらに事後 round lot が cutoff rule と整合的な銘柄だけを残す．
ここで outcome は Section 2 と同じ 7 本である．

\subsection{サンプル構成}

今回の RD 集計結果では，サンプル構成は次のとおりである．

\begin{center}
\begin{tabular}{lr}
\toprule
項目 & 値 \\
\midrule
分析対象の行数 & 118 \\
cutoff 左側の観測数 & 110 \\
cutoff 右側の観測数 & 8 \\
右側で実際に 40 へ移行した銘柄数 & 8 \\
cutoff & 250.0 \\
価格上限 & 1000.0 \\
\bottomrule
\end{tabular}
\end{center}

右側の観測数が 8 と小さい点は，この RD 分析の主要な制約である．

\subsection{RD 推定結果}

robust 標準誤差を使った RD 推定結果は次のとおりである．

\begin{table}[htbp]
\centering
\small
\resizebox{\linewidth}{!}{
\begin{tabular}{l l r r r r r r}
\toprule
アウトカム & 変化の定義 & 推定された跳び & 95\% CI 下限 & 95\% CI 上限 & p値 & 左側観測数 & 右側観測数 \\
\midrule
\texttt{cross\_k\_neg\_1e3\_1e4} & \( \log(1+Y^{\mathrm{post}})-\log(1+Y^{\mathrm{pre}}) \) & -0.418469 & -0.859487 &  0.022550 & 0.062921 & 110 & 8 \\
\texttt{cross\_k\_neg\_1e4\_1e5} & \( \log(1+Y^{\mathrm{post}})-\log(1+Y^{\mathrm{pre}}) \) & -0.215612 & -0.388886 & -0.042338 & 0.014733 & 110 & 8 \\
\texttt{cross\_k\_neg\_1e5\_0}   & \( \log(1+Y^{\mathrm{post}})-\log(1+Y^{\mathrm{pre}}) \) & -0.063431 & -0.167633 &  0.040771 & 0.232834 & 110 & 8 \\
\texttt{cross\_k\_pos\_0\_1e5}   & \( \log(1+Y^{\mathrm{post}})-\log(1+Y^{\mathrm{pre}}) \) & -0.070373 & -0.160050 &  0.019305 & 0.124039 & 110 & 8 \\
\texttt{cross\_k\_pos\_1e4\_1e3} & \( \log(1+Y^{\mathrm{post}})-\log(1+Y^{\mathrm{pre}}) \) & -0.609332 & -1.174146 & -0.044517 & 0.034477 & 110 & 8 \\
\texttt{cross\_k\_pos\_1e5\_1e4} & \( \log(1+Y^{\mathrm{post}})-\log(1+Y^{\mathrm{pre}}) \) & -0.252249 & -0.463223 & -0.041274 & 0.019109 & 110 & 8 \\
\texttt{mode}                    & \( Y^{\mathrm{post}}-Y^{\mathrm{pre}} \)                &  0.000017 & -0.000041 &  0.000076 & 0.558408 & 110 & 8 \\
\bottomrule
\end{tabular}
}
\caption{\$250\$ cutoff RD の推定結果．}
\end{table}

5\% 水準で負の discontinuity が確認されるのは
\texttt{cross\_k\_neg\_1e4\_1e5},
\texttt{cross\_k\_pos\_1e5\_1e4},
\texttt{cross\_k\_pos\_1e4\_1e3} の 3 つである．
10\% 水準では \texttt{cross\_k\_neg\_1e3\_1e4} も追加される．
\texttt{mode} の discontinuity は有意でない．

\subsection{RD 図}

\begin{figure}[p]
\centering
\includegraphics[width=0.78\linewidth]{rd_plot__mode.png}
\caption{\texttt{mode} の RD 図．}
\end{figure}

\begin{figure}[p]
\centering
\includegraphics[width=0.32\linewidth]{rd_plot__cross_k_neg_1e3_1e4.png}
\includegraphics[width=0.32\linewidth]{rd_plot__cross_k_neg_1e4_1e5.png}
\includegraphics[width=0.32\linewidth]{rd_plot__cross_k_neg_1e5_0.png}
\caption{負側の cross-K 変化の RD 図．左から
\texttt{cross\_k\_neg\_1e3\_1e4},
\texttt{cross\_k\_neg\_1e4\_1e5},
\texttt{cross\_k\_neg\_1e5\_0}.}
\end{figure}

\begin{figure}[p]
\centering
\includegraphics[width=0.32\linewidth]{rd_plot__cross_k_pos_0_1e5.png}
\includegraphics[width=0.32\linewidth]{rd_plot__cross_k_pos_1e5_1e4.png}
\includegraphics[width=0.32\linewidth]{rd_plot__cross_k_pos_1e4_1e3.png}
\caption{正側の cross-K 変化の RD 図．左から
\texttt{cross\_k\_pos\_0\_1e5},
\texttt{cross\_k\_pos\_1e5\_1e4},
\texttt{cross\_k\_pos\_1e4\_1e3}.}
\end{figure}

\subsection{制約}

\begin{itemize}
  \item \texttt{rddensity} は pandas 3 非互換のため omitted である．
  \item bandwidth source は \texttt{manual\_fallback} である．
  \item automatic bandwidth は
        \texttt{LinAlgError("Matrix is not positive definite")} により失敗している．
  \item Appendix E の balance table が示すように，
        事前 covariate にも強い discontinuity があるため，
        因果解釈は慎重であるべきである．
\end{itemize}

\clearpage
\appendix

\section{処理の数理的定義}

\subsection{event 列の定義}

trade 側では，日付 \(d\)，銘柄 \(s\)，取引所 \(e\) に対して，
観測された event time 列を
\[
  T_{d,s,e} = \{t^{\mathrm{tr}}_{d,s,e,i}\}_{i=1}^{N^{\mathrm{tr}}_{d,s,e}}
\]
と書く．ここで \(N^{\mathrm{tr}}_{d,s,e}\) は dedup 後の event 数であり，
\(i=1,\dots,N^{\mathrm{tr}}_{d,s,e}\) である．
各 \(t^{\mathrm{tr}}_{d,s,e,i}\) は \texttt{Participant Timestamp} を
秒へ変換した値であり，
\((\texttt{Exchange}, \texttt{Participant Timestamp})\) の重複を潰した後で
\(\mathcal{W}\) 内に制限したものとする．

quote 側では，日付 \(d\)，銘柄 \(s\)，取引所 \(g\) に対して，
\[
  Q_{d,s,g} = \{t^{\mathrm{qt}}_{d,s,g,j}\}_{j=1}^{N^{\mathrm{qt}}_{d,s,g}}
\]
を定義する．ここで \(N^{\mathrm{qt}}_{d,s,g}\) は dedup 後の event 数であり，
\(j=1,\dots,N^{\mathrm{qt}}_{d,s,g}\) である．
各 \(t^{\mathrm{qt}}_{d,s,g,j}\) は \texttt{Participant\_Timestamp} を
秒へ変換した値であり，
\((\texttt{Exchange}, \texttt{Participant\_Timestamp})\) の重複を潰した後で
\(\mathcal{W}\) 内に制限したものとする．

raw BBO から quote 側イベントを作る前に，
\[
  \texttt{Bid\_Price} > 0,\quad
  \texttt{Offer\_Price} > 0,\quad
  \texttt{Bid\_Price} < \texttt{Offer\_Price}
\]
を課す．

\subsection{pair summary の定義}

候補 pair 集合を
\[
  \mathcal{P}_{d,s}
  =
  \left\{
    (e,g)
    \mid
    e \in \mathcal{E}^{\mathrm{tr}}_{d,s},
    \ g \in \mathcal{E}^{\mathrm{qt}}_{d,s},
    \ e \neq g
  \right\}
\]
と書く．
pair \((e,g)\) は
\[
  N^{\mathrm{tr}}_{d,s,e} \ge 100,\quad
  N^{\mathrm{qt}}_{d,s,g} \ge 100
\]
を満たすときに解析対象になる．

共有 event time の個数
\[
  H_{d,s,e,g}
  =
  \#\left(T_{d,s,e}\cap Q_{d,s,g}\right)
\]
が正ならその pair は除外する．

解析対象の pair に対して，mode 集合
\[
  \mathcal{M}_{d,s,e,g}
  =
  \{m_{d,s,e,g,\ell}\}_{\ell=1}^{L_{d,s,e,g}}
\]
を求め，代表値
\[
  m^{\star}_{d,s,e,g}
  =
  \arg\min_{m \in \mathcal{M}_{d,s,e,g}} |m|
\]
を \texttt{closest\_mode\_to\_zero\_sec} として保持する．

cross-K は次の 6 区間で評価する．
\[
  [-10^{-3},-10^{-4}],
  [-10^{-4},-10^{-5}],
  [-10^{-5},0],
  [0,10^{-5}],
  [10^{-5},10^{-4}],
  [10^{-4},10^{-3}].
\]

\section{完全再現手順}

\subsection{前提}

\begin{itemize}
  \item repository root からコマンドを実行する．
  \item Python 環境は \texttt{uv} で同期する．
  \item Daily TAQ の MASTER, TRADE, BBO データは
        \texttt{data/dailyTAQ/} 配下にある前提で進める．
  \item 平均終値の付与にはネットワーク経由で \texttt{yfinance} を使う．
\end{itemize}

\subsection{panel 構築}

\begin{enumerate}
  \item 依存関係を同期する．
\begin{verbatim}
uv sync
\end{verbatim}

  \item M 接頭辞銘柄について filtered parquet を作る．
\begin{verbatim}
uv run python src/write_m_trade_quote_parquets.py
\end{verbatim}

  \item mode summary を計算する．
\begin{verbatim}
uv run python src/write_m_trade_quote_mode_summary.py \
  --date 20251031 \
  --date 20251103 \
  --symbol-prefix M \
  --min-events-per-side 100 \
  --n-jobs 8
\end{verbatim}

  \item \texttt{notebooks/m\_trade\_quote\_mode\_symbol\_panel.py} を開き，
        \texttt{run\_id} に直前の run directory 名を明示入力して全セル実行する．
\end{enumerate}

\subsection{round-lot change 回帰}

\begin{enumerate}
  \item \texttt{notebooks/m\_trade\_quote\_round\_lot\_effects.py} を開く．
  \item 既定の panel CSV を入力として全セル実行する．
  \item 出力先は
\begin{verbatim}
data/derived/m_trade_quote_round_lot_effects/
\end{verbatim}
        配下である．
\end{enumerate}

\subsection{RD 分析}

\begin{enumerate}
  \item \texttt{notebooks/m\_trade\_quote\_rd\_250.py} を開く．
  \item 既定の panel CSV を入力として全セル実行する．
  \item 出力先は
\begin{verbatim}
data/derived/m_trade_quote_rd_250/
\end{verbatim}
        配下である．
\end{enumerate}

\section{sanity check}

\begin{center}
\begin{tabular}{lll}
\toprule
対象 & 指標 & 期待値 \\
\midrule
MASTER 抽出 & 各日の M 接頭辞銘柄数 & 636, 636 \\
候補 pair & \(\#\mathcal{P}_{d,s}\) の総和 & 71,532; 73,555 \\
eligible pair & 両 side が 100 以上の pair 数 & 23,202; 24,053 \\
Q-Z ok pair & \texttt{status=ok} の行数 & 140; 143 \\
最終 panel & 行列サイズ & \((153,18)\) \\
最終 panel & 2 日とも値が揃う銘柄数 & 130 \\
最終 panel & round-lot changed symbols & 9 \\
\bottomrule
\end{tabular}
\end{center}

\section{回帰の群別要約}

\small
\begin{longtable}{p{0.22\linewidth}p{0.10\linewidth}r r r r r r}
\toprule
アウトカム & 群 & 観測数 & 平均価格 & 事前平均 & 事後平均 & 平均変化 & \(250\)ドル整合率 \\
\midrule
\endhead
\texttt{cross\_k\_neg\_1e3\_1e4} & 100のまま & 113 & 72.381 & 4.871217 & 4.715359 & -0.022200 & 0.973451 \\
\texttt{cross\_k\_neg\_1e3\_1e4} & 40へ移行  &   8 & 497.864 & 1.388735 & 1.346083 & -0.020540 & 1.000000 \\
\texttt{cross\_k\_neg\_1e4\_1e5} & 100のまま & 113 & 72.381 & 1.836563 & 1.745616 & -0.011688 & 0.973451 \\
\texttt{cross\_k\_neg\_1e4\_1e5} & 40へ移行  &   8 & 497.864 & 0.407685 & 0.372571 & -0.021090 & 1.000000 \\
\texttt{cross\_k\_neg\_1e5\_0}   & 100のまま & 113 & 72.381 & 0.183715 & 0.188516 &  0.003853 & 0.973451 \\
\texttt{cross\_k\_neg\_1e5\_0}   & 40へ移行  &   8 & 497.864 & 0.043183 & 0.038812 & -0.004045 & 1.000000 \\
\texttt{cross\_k\_pos\_0\_1e5}   & 100のまま & 113 & 72.381 & 0.159651 & 0.185225 &  0.014656 & 0.973451 \\
\texttt{cross\_k\_pos\_0\_1e5}   & 40へ移行  &   8 & 497.864 & 0.041407 & 0.038607 & -0.002621 & 1.000000 \\
\texttt{cross\_k\_pos\_1e4\_1e3} & 100のまま & 113 & 72.381 & 11.887236 & 11.489257 &  0.017207 & 0.973451 \\
\texttt{cross\_k\_pos\_1e4\_1e3} & 40へ移行  &   8 & 497.864 & 3.915405 & 3.428551 & -0.081487 & 1.000000 \\
\texttt{cross\_k\_pos\_1e5\_1e4} & 100のまま & 113 & 72.381 & 1.488521 & 1.534846 &  0.029585 & 0.973451 \\
\texttt{cross\_k\_pos\_1e5\_1e4} & 40へ移行  &   8 & 497.864 & 0.386890 & 0.388630 & -0.003616 & 1.000000 \\
\texttt{mode}                    & 100のまま & 113 & 72.381 & 0.000141 & 0.000142 &  0.000001 & 0.973451 \\
\texttt{mode}                    & 40へ移行  &   8 & 497.864 & 0.000159 & 0.000155 & -0.000004 & 1.000000 \\
\bottomrule
\end{longtable}
\normalsize

\section{RD balance test 詳細}

各 outcome について，pre-reform covariates として，
当該 outcome 自身と pre round lot を除く
\texttt{20251031} 列を balance test に使っている．

\small
\begin{longtable}{p{0.18\linewidth}p{0.28\linewidth}r r r r}
\toprule
outcome & covariate & 推定値 & p値 & 95\% CI 下限 & 95\% CI 上限 \\
\midrule
\endhead
\texttt{cross\_k\_neg\_1e3\_1e4} & \texttt{cross\_k\_pos\_1e4\_1e3\_20251031} & -8.307842 & 0.000014 & -12.061481 & -4.554203 \\
\texttt{cross\_k\_neg\_1e3\_1e4} & \texttt{cross\_k\_neg\_1e4\_1e5\_20251031} & -0.806629 & 0.000061 & -1.201112 & -0.412145 \\
\texttt{cross\_k\_neg\_1e3\_1e4} & \texttt{cross\_k\_neg\_1e5\_0\_20251031}   & -0.104168 & 0.000074 & -0.155691 & -0.052646 \\
\texttt{cross\_k\_neg\_1e3\_1e4} & \texttt{cross\_k\_pos\_1e5\_1e4\_20251031} & -0.867424 & 0.000080 & -1.298330 & -0.436518 \\
\texttt{cross\_k\_neg\_1e3\_1e4} & \texttt{cross\_k\_pos\_0\_1e5\_20251031}   & -0.062949 & 0.015566 & -0.113954 & -0.011944 \\
\texttt{cross\_k\_neg\_1e3\_1e4} & \texttt{mode\_20251031}                    & -0.000027 & 0.204805 & -0.000069 &  0.000015 \\
\texttt{cross\_k\_neg\_1e4\_1e5} & \texttt{cross\_k\_pos\_1e4\_1e3\_20251031} & -8.307842 & 0.000014 & -12.061481 & -4.554203 \\
\texttt{cross\_k\_neg\_1e4\_1e5} & \texttt{cross\_k\_neg\_1e3\_1e4\_20251031} & -3.794987 & 0.000060 & -5.648330 & -1.941645 \\
\texttt{cross\_k\_neg\_1e4\_1e5} & \texttt{cross\_k\_neg\_1e5\_0\_20251031}   & -0.104168 & 0.000074 & -0.155691 & -0.052646 \\
\texttt{cross\_k\_neg\_1e4\_1e5} & \texttt{cross\_k\_pos\_1e5\_1e4\_20251031} & -0.867424 & 0.000080 & -1.298330 & -0.436518 \\
\texttt{cross\_k\_neg\_1e4\_1e5} & \texttt{cross\_k\_pos\_0\_1e5\_20251031}   & -0.062949 & 0.015566 & -0.113954 & -0.011944 \\
\texttt{cross\_k\_neg\_1e4\_1e5} & \texttt{mode\_20251031}                    & -0.000027 & 0.204805 & -0.000069 &  0.000015 \\
\texttt{cross\_k\_neg\_1e5\_0}   & \texttt{cross\_k\_pos\_1e4\_1e3\_20251031} & -8.307842 & 0.000014 & -12.061481 & -4.554203 \\
\texttt{cross\_k\_neg\_1e5\_0}   & \texttt{cross\_k\_neg\_1e3\_1e4\_20251031} & -3.794987 & 0.000060 & -5.648330 & -1.941645 \\
\texttt{cross\_k\_neg\_1e5\_0}   & \texttt{cross\_k\_neg\_1e4\_1e5\_20251031} & -0.806629 & 0.000061 & -1.201112 & -0.412145 \\
\texttt{cross\_k\_neg\_1e5\_0}   & \texttt{cross\_k\_pos\_1e5\_1e4\_20251031} & -0.867424 & 0.000080 & -1.298330 & -0.436518 \\
\texttt{cross\_k\_neg\_1e5\_0}   & \texttt{cross\_k\_pos\_0\_1e5\_20251031}   & -0.062949 & 0.015566 & -0.113954 & -0.011944 \\
\texttt{cross\_k\_neg\_1e5\_0}   & \texttt{mode\_20251031}                    & -0.000027 & 0.204805 & -0.000069 &  0.000015 \\
\texttt{cross\_k\_pos\_0\_1e5}   & \texttt{cross\_k\_pos\_1e4\_1e3\_20251031} & -8.307842 & 0.000014 & -12.061481 & -4.554203 \\
\texttt{cross\_k\_pos\_0\_1e5}   & \texttt{cross\_k\_neg\_1e3\_1e4\_20251031} & -3.794987 & 0.000060 & -5.648330 & -1.941645 \\
\texttt{cross\_k\_pos\_0\_1e5}   & \texttt{cross\_k\_neg\_1e4\_1e5\_20251031} & -0.806629 & 0.000061 & -1.201112 & -0.412145 \\
\texttt{cross\_k\_pos\_0\_1e5}   & \texttt{cross\_k\_neg\_1e5\_0\_20251031}   & -0.104168 & 0.000074 & -0.155691 & -0.052646 \\
\texttt{cross\_k\_pos\_0\_1e5}   & \texttt{cross\_k\_pos\_1e5\_1e4\_20251031} & -0.867424 & 0.000080 & -1.298330 & -0.436518 \\
\texttt{cross\_k\_pos\_0\_1e5}   & \texttt{mode\_20251031}                    & -0.000027 & 0.204805 & -0.000069 &  0.000015 \\
\texttt{cross\_k\_pos\_1e4\_1e3} & \texttt{cross\_k\_neg\_1e3\_1e4\_20251031} & -3.794987 & 0.000060 & -5.648330 & -1.941645 \\
\texttt{cross\_k\_pos\_1e4\_1e3} & \texttt{cross\_k\_neg\_1e4\_1e5\_20251031} & -0.806629 & 0.000061 & -1.201112 & -0.412145 \\
\texttt{cross\_k\_pos\_1e4\_1e3} & \texttt{cross\_k\_neg\_1e5\_0\_20251031}   & -0.104168 & 0.000074 & -0.155691 & -0.052646 \\
\texttt{cross\_k\_pos\_1e4\_1e3} & \texttt{cross\_k\_pos\_1e5\_1e4\_20251031} & -0.867424 & 0.000080 & -1.298330 & -0.436518 \\
\texttt{cross\_k\_pos\_1e4\_1e3} & \texttt{cross\_k\_pos\_0\_1e5\_20251031}   & -0.062949 & 0.015566 & -0.113954 & -0.011944 \\
\texttt{cross\_k\_pos\_1e4\_1e3} & \texttt{mode\_20251031}                    & -0.000027 & 0.204805 & -0.000069 &  0.000015 \\
\texttt{cross\_k\_pos\_1e5\_1e4} & \texttt{cross\_k\_pos\_1e4\_1e3\_20251031} & -8.307842 & 0.000014 & -12.061481 & -4.554203 \\
\texttt{cross\_k\_pos\_1e5\_1e4} & \texttt{cross\_k\_neg\_1e3\_1e4\_20251031} & -3.794987 & 0.000060 & -5.648330 & -1.941645 \\
\texttt{cross\_k\_pos\_1e5\_1e4} & \texttt{cross\_k\_neg\_1e4\_1e5\_20251031} & -0.806629 & 0.000061 & -1.201112 & -0.412145 \\
\texttt{cross\_k\_pos\_1e5\_1e4} & \texttt{cross\_k\_neg\_1e5\_0\_20251031}   & -0.104168 & 0.000074 & -0.155691 & -0.052646 \\
\texttt{cross\_k\_pos\_1e5\_1e4} & \texttt{cross\_k\_pos\_0\_1e5\_20251031}   & -0.062949 & 0.015566 & -0.113954 & -0.011944 \\
\texttt{cross\_k\_pos\_1e5\_1e4} & \texttt{mode\_20251031}                    & -0.000027 & 0.204805 & -0.000069 &  0.000015 \\
\texttt{mode}                    & \texttt{cross\_k\_pos\_1e4\_1e3\_20251031} & -8.307842 & 0.000014 & -12.061481 & -4.554203 \\
\texttt{mode}                    & \texttt{cross\_k\_neg\_1e3\_1e4\_20251031} & -3.794987 & 0.000060 & -5.648330 & -1.941645 \\
\texttt{mode}                    & \texttt{cross\_k\_neg\_1e4\_1e5\_20251031} & -0.806629 & 0.000061 & -1.201112 & -0.412145 \\
\texttt{mode}                    & \texttt{cross\_k\_neg\_1e5\_0\_20251031}   & -0.104168 & 0.000074 & -0.155691 & -0.052646 \\
\texttt{mode}                    & \texttt{cross\_k\_pos\_1e5\_1e4\_20251031} & -0.867424 & 0.000080 & -1.298330 & -0.436518 \\
\texttt{mode}                    & \texttt{cross\_k\_pos\_0\_1e5\_20251031}   & -0.062949 & 0.015566 & -0.113954 & -0.011944 \\
\bottomrule
\end{longtable}
\normalsize

\section{再現上の注意点}

\begin{itemize}
  \item 日次終値は \texttt{yfinance} 依存なので，ネットワーク状況や取得失敗により
        \(\bar{p}^{\mathrm{close}}_s\) の欠損が増える可能性がある．
  \item panel notebook では \texttt{run\_id} を空欄にすると latest directory を
        自動採用するため，共同研究者へ共有する際は必ず明示的に run 名を入れる．
  \item RD の既存成果物では dirty worktree が前提になっているので，
        完全一致よりも sample size，係数，信頼区間の整合をまず確認するべきである．
  \item \texttt{notebooks/m\_trade\_quote\_round\_lot\_effects.py} と
        \texttt{notebooks/m\_trade\_quote\_rd\_250.py} は現時点では working tree 上の
        notebook を前提としている．
\end{itemize}

\end{document}
